\section{Lean formalization architecture}
\label{sec:lean-plan}

\subsection{Target environment}

The formalization should live inside the existing ProveIt Lean workspace and
import the checked modules rather than copying a second kernel:

\begin{lstlisting}[style=lean]
import ConciseIPRA.PR
import ConciseIPRA.NativeBiInterpretation

namespace LeanProofs.ConciseIPRA.Finitary
open Syn
open scoped Fm
open PTm (app zero succ)
\end{lstlisting}

The source-accurate interfaces used by the plan are:

\begin{itemize}
\item \code{PR}, \code{PTm}, and \code{PRAAx};
\item \code{PRA.precZero}, \code{PRA.precSucc}, and \code{PRA.qfInd};
\item the constructors and structural lemmas of \code{Prf}, including
      substitution, weakening, cut, universal introduction/elimination, and
      equality substitution; and
\item \code{pra\_iff\_native} for final transport.
\end{itemize}

The development should not start with Mathlib real numbers.  Ambient finite
models are useful for testing, but the primary endpoint of each generic lemma
is a value of type \code{Prf PRAAx Gamma phi}.

\subsection{Canonical code library}

The first module, proposed as \code{Finitary/Codes.lean}, chooses exactly one
code for each recurring data type:

\begin{enumerate}
\item Booleans normalized to zero and one;
\item pairs and bounded tuples;
\item bounded lists and bit vectors;
\item signed integers;
\item positive-denominator rationals and dyadics;
\item finite maps represented by bounded tables.
\end{enumerate}

For every operation there are three layers:

\begin{enumerate}
\item a \code{PR} symbol;
\item a \code{PTm} notation/helper for applying the symbol; and
\item embedded defining/correctness derivations.
\end{enumerate}

Parallel encodings should be avoided.  Most proof-packet integration failures
would be amplified if each packet introduced its own rational or list code.

\subsection{Boolean folds and bounded search}

The next modules are \code{Boolean.lean}, \code{Fold.lean}, and
\code{Search.lean}.  Suggested theorem signatures are schematic generators:

\begin{lstlisting}[style=lean]
/-- Object-language proof that a fixed PR Boolean fold is correct. -/
def allFoldCorrect (b : PR (r + 1)) :
  [] |-[PRAAx] allFoldSpec b := ...

def anyFoldCorrect (b : PR (r + 1)) :
  [] |-[PRAAx] anyFoldSpec b := ...

def leastSearchCorrect (b : PR (r + 1)) :
  [] |-[PRAAx] leastSearchSpec b := ...
\end{lstlisting}

The formulas \code{allFoldSpec}, \code{anyFoldSpec}, and
\code{leastSearchSpec} are built from equations and bounded numerical
relations, not Lean propositions about function values.  Their proofs use
\code{PRA.qfInd}; the logical bounded-quantifier corollaries are derived after
the Boolean equations.

\subsection{Trace-system definitions}

The fixed symbols of a trace are represented directly:

\begin{lstlisting}[style=lean]
structure PRTraceSystem where
  good   : PR 1
  init   : PR 1
  bound  : PR 1
  step   : PR 3      -- (i, previous, u)
  inv    : PR 3      -- (i, state, u)
  accept : PR 2      -- (state, u)

def PRTraceSystem.run (T : PRTraceSystem) : PR 2 :=
  .prec T.init T.step
\end{lstlisting}

Application helpers for arities one through three should be defined using
\code{Fin.cases}.  Define \code{IsOne t := t = succ zero}.  With the de Bruijn
convention documented in the source, define open formulas:

\begin{lstlisting}[style=lean]
def BaseFm  (T : PRTraceSystem) : Fm PTm 1 := ... -- free u
def StepFm  (T : PRTraceSystem) : Fm PTm 3 := ... -- free i,s,u
def FinalFm (T : PRTraceSystem) : Fm PTm 2 := ... -- free s,u
def ResultFm (T : PRTraceSystem) : Fm PTm 0 :=
  Fm.all (... Good(u) -> Accept(run(bound u,u),u) ...)

structure TraceProofs (T : PRTraceSystem) where
  base  : [] |-[PRAAx] BaseFm T
  step  : [] |-[PRAAx] StepFm T
  final : [] |-[PRAAx] FinalFm T
\end{lstlisting}

The intended central theorem is:

\begin{lstlisting}[style=lean]
def traceReflection (T : PRTraceSystem) (h : TraceProofs T) :
    [] |-[PRAAx] ResultFm T := by
  -- instantiate h.base at u
  -- rewrite run(0,u) using PRA.precZero
  -- apply PRA.qfInd to Inv(i, run(i,u), u)
  -- in the step, instantiate h.step and rewrite with PRA.precSucc
  -- specialize the universal invariant at bound(u)
  -- instantiate h.final and cut
  -- discharge Good(u), then universally introduce u
  ...
\end{lstlisting}

This theorem is deliberately a \code{def}, not merely an existence theorem:
it computes the final derivation from the three premise derivations.  The
construction should use only the existing proof constructors, so its axiom
audit is inherited from \code{ConciseIPRA}.

\subsection{Proof-producing finite-program DSL}

After the raw trace theorem works, define a small typed syntax of finite
programs:

\begin{lstlisting}[style=lean]
inductive FProg (input output : Nat) where
  | pure    : PR input -> FProg input output
  | seq     : FProg a b -> FProg b c -> FProg a c
  | loop    : PR input -> PR state -> PR (...) -> FProg input output
  | branch  : PR (...) -> FProg input output ->
              FProg input output -> FProg input output
  | check   : PR (...) -> FProg input input
\end{lstlisting}

The actual indices should track state arities and contracts more carefully than
this sketch.  Compile every program to a \code{PRTraceSystem}.  A
\code{CertifiedProg} pairs syntax with local \code{Prf PRAAx} obligations.
Prove by structural recursion that compilation returns the global soundness
proof.  This automates phase composition and eliminates repeated hand-written
induction scaffolding across theorem packets.

\subsection{Relational finite kernel}

A semantic prototype can use ordinary Lean finite data, but the integrated
module should expose fixed numerical symbols:

\begin{lstlisting}[style=lean]
structure PRLayeredRelation where
  bound  : PR 2       -- (k,p)
  good   : PR 3       -- (a,k,p)
  compat : PR 4       -- (b,a,k,p)

def viable (R : PRLayeredRelation) : PR 4 := ...
def reach (R : PRLayeredRelation) : PR 3 := ...
def fullViable (R : PRLayeredRelation) : PR 4 := ...
\end{lstlisting}

Suggested embedded endpoints are:

\begin{lstlisting}[style=lean]
def viableIffPath (R : PRLayeredRelation) :
  [] |-[PRAAx] viablePathSpec R := ...

def leastViablePathSound (R : PRLayeredRelation) :
  [] |-[PRAAx] leastPathSpec R := ...

def fullViableSplit (R : PRLayeredRelation) :
  [] |-[PRAAx] fullViableSpec R := ...
\end{lstlisting}

Implement viability first as a depth recursion, then add the table algorithm
and prove the two implementations equal.  The depth recursion has the simplest
proof; the table version has the useful execution complexity.

\subsection{Shrinking-viability module}

The module \code{Finitary/Shrinking.lean} should parameterize over coded
rational observations and a coded pseudometric.  All assumptions are
object-language derivations:

\begin{lstlisting}[style=lean]
structure ShrinkingProofs (R : PRLayeredRelation)
    (obs : PR 3) (dist eps delta horizon : ...) where
  nonempty : [] |-[PRAAx] NonemptyCloudFm ...
  monotone : [] |-[PRAAx] HorizonMonotoneFm ...
  diameter : [] |-[PRAAx] CloudDiameterFm ...
  edge     : [] |-[PRAAx] EdgeBoundFm ...
  metric   : MetricProofs dist

def shrinkingViability (h : ShrinkingProofs ...) :
  [] |-[PRAAx] ShrinkingBoundFm ... := ...
\end{lstlisting}

The proof dependencies are:

\begin{enumerate}
\item least bounded search in the full-viable cloud;
\item extraction of a least path certificate from \code{fullViable};
\item truncation of a valid path;
\item finite telescoping of edge distances; and
\item one triangle inequality combining cloud diameter and path motion.
\end{enumerate}

A corollary consumes a tail modulus and constructs a regular-name theorem.
Do not define a quotient real or invoke classical choice.

\subsection{Rational and precision modules}

The modules \code{Rational.lean}, \code{ErrorBudget.lean}, and
\code{RegularName.lean} should prove:

\begin{itemize}
\item cross-multiplication correctness under positive denominators;
\item triangle and product-error inequalities;
\item finite-sum and finite-product accumulators;
\item regularity/extensionality of addition and multiplication;
\item totalized inverse under an apartness guard;
\item overlap branching and exact branching under a gap; and
\item composition of precision ledgers.
\end{itemize}

Most proofs should normalize to polynomial integer inequalities and use small
reusable order lemmas.  A reflection tactic may generate proof terms for
closed numeral calculations, but the universal algebraic lemmas must remain
explicit derivations in the object theory.

\subsection{Application interface}

Each theorem packet should expose a record of the following form:

\begin{lstlisting}[style=lean]
structure PacketTarget where
  params : Nat
  assumptions : List (Fm PTm params)
  conclusion  : Fm PTm params
  derivation  : assumptions |-[PRAAx] conclusion

 def PacketTarget.toNative (T : PacketTarget) :
   NPrf (T.assumptions.map NFm.ofPRA) (NFm.ofPRA T.conclusion) :=
   (pra_iff_native).mp T.derivation
\end{lstlisting}

For an incomplete packet, replace \code{derivation} by an explicitly named
\code{Goal} proposition or a structure of missing producer obligations.  Do
not fill the field with an axiom, \code{opaque} placeholder, or proof of an
ambient semantic surrogate.

\subsection{Dependency-ordered milestones}

A realistic implementation order is:

\begin{enumerate}
\item application helpers and Boolean truth;
\item pairing, bounded lists, signed integers, rationals;
\item folds, bounded minimization, path codes;
\item embedded trace reflection;
\item compiler composition and finite-program DSL;
\item reachability, viability, and least paths;
\item rational metric and shrinking viability;
\item regular names, guards, robust branching, finite nets;
\item exact polynomial/matrix kernel;
\item domain adapters and packet migrations;
\item native export and proof-size measurements.
\end{enumerate}

The first vertical slice should be deliberately small: formalize a finite sum
trace, produce its exact \(\Prf\,\PRAAx\) theorem, and transport it to
\(\NPrf\).  The second slice should be the square-root interval example from
\cref{sec:worked-examples}, because it exercises relational candidates,
viability, shrinking, and rational error bounds.

\subsection{Audit requirements}

For every completed endpoint:

\begin{enumerate}
\item run the repository Lake build;
\item scan source for \code{sorry}, \code{admit}, and new \code{axiom}
      declarations;
\item use \code{\#print axioms} on public theorems;
\item distinguish definitions of goals from proofs of goals;
\item record the pinned ProveIt commit and Lean toolchain;
\item preserve a theorem ledger mapping human claims to Lean names; and
\item retain exact executable checks as regression tests, not as replacements
      for universal proofs.
\end{enumerate}

The accompanying Lean files in this archive are intentionally blueprints: no
Lean executable is available in the build environment, and the generic
\(\Prf\,\PRAAx\) compiler has not been kernel-checked here.  Their purpose is
to make the remaining implementation path explicit and source-accurate.
